\documentclass[11pt]{article}
\usepackage[margin=1in]{geometry}
\usepackage[T1]{fontenc}
\usepackage{amsmath,amssymb}
\usepackage{booktabs}
\usepackage{enumitem}
\usepackage{hyperref}
\usepackage{xcolor}

\newcommand{\code}[1]{\texttt{\detokenize{#1}}}
\pdfstringdefDisableCommands{\def\code#1{#1}}
\newcommand{\gradr}{\langle|\nabla r|^2\rangle}
\newcommand{\Dped}{D_{\mathrm{ped}}}
\newcommand{\nFC}{\langle n_{\mathrm{FC}}\rangle}
\newcommand{\nCX}{\langle n_{\mathrm{CX}}\rangle}
\newcommand{\SCX}{S_{\mathrm{CX}}}
\newcommand{\VFC}{|V_{\mathrm{FC}}|}
\newcommand{\VCX}{|V_{\mathrm{CX}}|}
\newcommand{\fFC}{f_{\mathrm{FC}}}
\newcommand{\fCX}{f_{\mathrm{CX}}}
\newcommand{\CCX}{C_{\mathrm{CX}}}

\title{\code{src/solver_1d.py}: Documentation of the Original (Implicit)\\
Saarelma--Connor Pedestal Solver (``1D'' model)}
\author{}
\date{}

\begin{document}
\maketitle
\tableofcontents

%=====================================================================
\section{Overview}
%=====================================================================
\code{src/solver_1d.py} contains non-dimensional solvers for the
\emph{original} (implicit) Saarelma--Connor pedestal model, in both a
Firedrake and a scipy implementation. The module contributes the mixin
\code{OneDSolverMixin} to the single public class
\code{src.saarelma_connor_api.saarelma_connor}; all equilibrium, kinetic
and cross-section initialisation comes from \code{SaarelmaConnorBase}.

The model is the single implicit second-order ODE for the electron density
obtained after the two neutral fluids have been eliminated analytically
(S.~Saarelma \emph{et al.} 2023 Nucl.\ Fusion \textbf{63} 052002;
Eqs.~(6)--(7) of the Labbate APAM E9301 report; rigorous derivations in
\code{docs/derivation_eq16.tex} and \code{docs/derivation_eq15.tex}). Only
the electron density is solved for: the neutrals never appear as unknowns,
they enter through the closed-form closures already substituted into
Eqs.~(6)--(7).

\paragraph{Entry points.}
\begin{itemize}[nosep]
  \item \code{saarelma_connor.solve_sc_scipy(...)} --- scipy \code{solve_bvp};
  \item \code{saarelma_connor.solve_sc_firedrake(...)} --- Firedrake / SNES;
  \item \code{saarelma_connor.solve_sc(...)} --- thin dispatcher;
  \item \code{saarelma_connor.solve(model='1D', ...)} --- unified entry point.
\end{itemize}
All of them return the common result dictionary described in
\code{SaarelmaConnorBase._build_result_dict}.

\paragraph{Free parameters.} $\alpha_{\mathrm{crit}}$,
\code{De_chie_etg}, $C_{\mathrm{KBM}}$, \code{nFC_x0}; set through the
parent constructor, through \code{update_free_params}, or per solve through
the \code{free_params} argument of either entry point.

%=====================================================================
\section{Model}
%=====================================================================
\subsection{Step 1: the no-CX ``first step'' (report Eq.~6 $\equiv$ Saarelma Eq.~16)}

Two forms are available, selected by \code{eq6_form}:
\begin{align}
\frac{d}{dx}\!\left[\gradr \Dped \frac{dn_e}{dx}\right]
 &= \frac{n_e \Dped (S_i+\SCX)}{\VFC \fFC}
    \left(\frac{dn_e}{dx} - \left.\frac{dn_e}{dx}\right|_{\mathrm{in}}\right)
 && (\code{eq6_form="paper"}), \label{eq:eq6paper}\\
\frac{d}{dx}\!\left[\gradr \Dped \frac{dn_e}{dx}\right]
 &= \frac{n_e (S_i+\SCX)}{1+\SCX/(2S_i)}\,
    \frac{\gradr \Dped}{\VFC \fFC}
    \left(\frac{dn_e}{dx} - \left.\frac{dn_e}{dx}\right|_{\mathrm{in}}\right)
 && (\code{eq6_form="complete"}). \label{eq:eq6complete}
\end{align}

Eq.~\eqref{eq:eq6paper} is the equation exactly as printed;
Eq.~\eqref{eq:eq6complete} is the complete form of
\code{docs/derivation_eq16.tex} Eq.~(16-full). The printed form follows
from the complete one by dropping the shape factor $1+\SCX/(2S_i)$
(assumption A5) and by cancelling the $\gradr$ of the transport operator
against the $\Dped$ of the neutral closure --- a \emph{symbolic}
cancellation that is only legitimate when $\gradr\sim 1$.

The two differ by the factor $(1+\SCX/(2S_i))/\gradr$, which is not a
detail: on the DIII-D 158091 case shipped with this repository
$\gradr\sim 0.05$, so the printed form is $\sim 35\times$ stiffer than the
complete one. Its solution amplifies by $e^{176}$ across the pedestal and
neither implementation can solve it in double precision, whereas the
complete form amplifies by $\sim 70$ and converges. \code{"complete"} is
therefore the default; it is also the form consistent with Eq.~(7) below,
which keeps $\gradr\Dped$ in the same place.

\subsection{Step 2: the full implicit equation (report Eq.~7 $\equiv$ corrected Saarelma Eq.~15)}
\begin{equation}
\frac{d}{dx}\!\left[\gradr \Dped \frac{dn_e}{dx}\right]
 = -n_e S_i \left\{ -\frac{\gradr\Dped}{\VCX\fCX}
   \left(\frac{dn_e}{dx} - \left.\frac{dn_e}{dx}\right|_{\mathrm{in}}\right)
   + \CCX(x)\,\nFC(0)\,E(x) \right\},
\label{eq:eq7}
\end{equation}
with the CX-flux coefficient and the FC attenuation (``optical depth'')
kernel
\begin{align}
\CCX(x) &= 1 - \frac{\VFC\fFC}{\VCX\fCX}\,\frac{S_i+\SCX/2}{S_i+\SCX},
\label{eq:Ccx}\\
E(x) &= \exp\!\left[\int_0^x \frac{n_e(x')(S_i+\SCX)}{\fFC\VFC}\,dx'\right].
\label{eq:Ekernel}
\end{align}
The form factors $\fFC$ and $\fCX$ are carried explicitly. The report
writes Eq.~6 with $\fFC=1$ (assumption A4 of
\code{docs/derivation_eq16.tex}, which is what
\code{saarelma_connor.form_factor} sets), so the two forms agree
identically for the current form-factor model.

\subsection{Pedestal diffusivity}
\begin{equation}
\Dped(x,n_e) = D_{\mathrm{NEO}}(x) + D_{\mathrm{KBM}}(x) + \frac{C_{\mathrm{ETG}}(x)}{n_e},
\end{equation}
with $C_{\mathrm{ETG}} = \code{De_chie_etg}\cdot P_{\mathrm{tot},e}/(S_{\mathrm{plasma}}|dT_e/dx|)$
and $D_{\mathrm{KBM}}$ from the pedestal-averaged Connor--Hastie $\alpha$
(Saarelma Eqs.~24--25); see \code{calc_D_KBM_average_sc}
(Section~\ref{sec:kbm}).

%=====================================================================
\section{Implicit (Picard) solve}
%=====================================================================
Following Saarelma \emph{et al.} (2023), the non-local kernel $E(x)$ makes
Eq.~\eqref{eq:eq7} an integro-differential equation, so it is solved
implicitly:
\begin{enumerate}
  \item Solve Eq.~(6) (no CX neutrals) for a first $n_e$ profile.
  \item Iterate Eq.~(7): on each Picard iteration $E(x)$ is evaluated with
  the \emph{previous} $n_e$ and frozen, and $D_{\mathrm{KBM}}$ is refrozen
  from the latest $n_e$ using the pedestal-averaged (``average'' gate mode)
  $\alpha$ --- exactly the treatment of the 3D solver's
  \code{solve_coupled_nondim} with \code{kbm_treatment="picard"} /
  \code{picard_gate_mode="average"} (frozen scalar-gated diffusivity,
  optional under-relaxation). Repeat until the profile and the gate state
  stop changing.
\end{enumerate}

Two practical notes on this loop, both observed on the DIII-D 158091 case
and both controllable from the solver signature:
\begin{itemize}
  \item \textbf{Step 1 can have no solution.} Eq.~(6) is a positive-feedback
  amplifier (more density gradient $\to$ larger implied $\nFC$ $\to$ more
  ionisation $\to$ more gradient), and for a steep enough inner-boundary
  slope every profile collapses to $n_e\sim 0$ before reaching the
  separatrix, so no member of the family satisfies $n_e(0)=\code{ne_x0}$.
  Saarelma \emph{et al.} only use Eq.~(16) as ``a convenient starting
  point'', so \code{first_step} selects \code{"eq6"} (always run it, raise
  on failure), \code{"skip"} (go straight to the Eq.~7 Picard loop from
  \code{initial_guess}) or \code{"auto"} (default: try Eq.~6, warn and fall
  back to \code{"skip"}).
  \item \textbf{The KBM gate can limit-cycle.} When the converged
  $\bar\alpha$ sits close to $\alpha_{\mathrm{crit}}$ the whole-pedestal
  gate flips every iteration and the profile alternates between two states
  forever. Damp it with \code{picard_relax} $<1$ (0.5 works on the case
  above); the non-convergence error message diagnoses this explicitly.
\end{itemize}

%=====================================================================
\section{Non-dimensionalisation}
%=====================================================================
\begin{align*}
\xi &= x/L, & L &= |x_{\mathrm{inner}}|, & \xi&\in[-1,0],\\
N &= n_e/n_0, & n_0 &= \code{ne_x0}, & N(0)&=1,\\
S_0 &= S_i(0) \text{ (separatrix ionisation rate coefficient)}, &
[V]_0 &= L n_0 S_0, & [D]_0 &= L^2 n_0 S_0.
\end{align*}

Writing $f(x,n_e)=\gradr\Dped$ for the total conductance and splitting off
its $n_e$ dependence, $f = f_0(x) + f_1(x)/n_e$ with
\[
f_0 = \gradr(D_{\mathrm{NEO}}+D_{\mathrm{KBM}}),\qquad f_1 = \gradr C_{\mathrm{ETG}},
\]
the expanded non-dimensional ODEs solved by the \textbf{scipy}
implementation are
\begin{align}
\text{Eq.~6:}\quad N'' &= C_{A6}\,N(N'-N'_{\mathrm{in}}) - C_K N' + C_N (N')^2,\\
\text{Eq.~7:}\quad N'' &= C_A\,N(N'-N'_{\mathrm{in}}) - C_E N - C_K N' + C_N (N')^2,
\end{align}
with
\begin{align*}
C_{A6} &= \frac{n_0 L (S_i+\SCX)}{\gradr\VFC\fFC} \quad(\code{"paper"}),\qquad
C_{A6} = \frac{n_0 L (S_i+\SCX)}{(1+\SCX/(2S_i))\VFC\fFC} \quad(\code{"complete"}),\\
C_A &= \frac{n_0 L S_i}{\VCX\fCX},\qquad
C_E = \frac{L^2 S_i \CCX \nFC(0) E(x)}{f},\\
C_K &= \frac{L}{f}\left(\frac{df_0}{dx} + \frac{1}{n_e}\frac{df_1}{dx}\right),\qquad
C_N = \frac{f_1}{n_0 N^2 f},
\end{align*}
where $N'_{\mathrm{in}} = (L/n_0)\,dn_e/dx|_{\mathrm{in}}$ is the
non-dimensional Neumann value (cf.\ \code{docs/derivation_eq16.tex},
Sec.~``Non-dimensionalization'').

The \textbf{Firedrake} implementation keeps the conservative form. With
$\hat f=f/[D]_0$, $\hat S=S/S_0$, $\hat V=V/[V]_0$ and
$\hat n_{\mathrm{FC},0}=\code{nFC_x0}/n_0$ the strong forms are
\begin{align*}
\text{Eq.~6 (paper):}\quad & \frac{d}{d\xi}[\hat f N'] = N\,\frac{\hat f}{\gradr}\,
   \frac{(\hat S_i+\hat S_{\mathrm{CX}})(N'-N'_{\mathrm{in}})}{\hat V_{\mathrm{FC}}\fFC},\\
\text{Eq.~6 (complete):}\quad & \frac{d}{d\xi}[\hat f N'] = N\,\hat f\,
   \frac{\hat S_i+\hat S_{\mathrm{CX}}}{1+\hat S_{\mathrm{CX}}/(2\hat S_i)}\,
   \frac{N'-N'_{\mathrm{in}}}{\hat V_{\mathrm{FC}}\fFC},\\
\text{Eq.~7:}\quad & \frac{d}{d\xi}[\hat f N'] = N\hat S_i\left[
   \frac{\hat f(N'-N'_{\mathrm{in}})}{\hat V_{\mathrm{CX}}\fCX}
   - \CCX\,\hat n_{\mathrm{FC},0}\,E(\xi)\right],
\end{align*}
which are multiplied by a test function $w$ and integrated by parts in the
usual way (the residuals are assembled in \code{solve_sc_firedrake},
Section~\ref{sec:firedrake}).

\subsection{Agreement between the two implementations}
On the DIII-D 158091 case the converged profiles agree to
$\sim 3\times10^{-3}$ in max-norm, and each satisfies its own discrete
Eq.~(7) far better than that (the conservative residual of the Firedrake
solution falls as the mesh is refined, $\sim10^{-3}$ at \code{x_res = 400};
the scipy solution sits at $\sim 2\times10^{-4}$ of its own collocation
problem). The residual difference between them is dominated not by the
discretisation but by how each reconstructs the equilibrium coefficients
between the points of the (coarse) p-file grid: the Firedrake path samples
them with \code{np.interp} into the finite-element space, as the rest of
this repository does, while the scipy path needs $df/dx$ and therefore
uses a shape-preserving PCHIP interpolant.

%=====================================================================
\section{Boundary conditions}
\label{sec:bcs}
%=====================================================================
\code{ne_grad_bc_loc} selects which end carries the prescribed gradient.

\paragraph{\code{"inner"} (default, unchanged behaviour).}
\begin{align*}
\text{Neumann at } \xi=-1 \text{ (inner)}:&\quad N'(-1) = (L/n_0)\,dn_e/dx|_{\mathrm{in}},\\
\text{Dirichlet at } \xi=0 \text{ (outer)}:&\quad N(0) = 1 \quad(n_e(0)=\code{ne_x0}).
\end{align*}

\paragraph{\code{"outer"} --- both conditions at the separatrix, inner boundary free.}
\begin{align*}
\text{Dirichlet at } \xi=0:&\quad N(0) = 1,\\
\text{Neumann at } \xi=0:&\quad N'(0) = (L/n_0)\,dn_e/dx|_0,
\end{align*}
with nothing imposed at $\xi=-1$. This is the boundary-condition set of
Saarelma \emph{et al.} 2023 Sec.~2.3, where the separatrix density and its
gradient --- the latter from their SOL model, Eq.~(20):
$dn_e/dx|_0 = -n_e(0)/\sqrt{D_{\mathrm{SOL}}\tau_\parallel}$ --- are the
specified data and nothing is imposed at the pedestal top.

\paragraph{The two roles of $N'_{\mathrm{in}}$.} $N'_{\mathrm{in}}$ is
\emph{two} different things in this model and only one of them is a
boundary condition:
\begin{itemize}[nosep]
  \item the value of the Neumann condition at $\xi=-1$, and
  \item Saarelma's constant of integration $C = dn_e/dx|_{x=-\infty}$, which
  appears inside the source term of both Eq.~(6) and Eq.~(7) as the
  combination $(N'-N'_{\mathrm{in}})$ --- the accumulated ionisation source
  between the deep interior and $x$, i.e.\ the statement that the neutrals
  are extinguished by the pedestal top.
\end{itemize}
In \code{"inner"} mode the two coincide, which is self-consistent: imposing
$N'(-1)=N'_{\mathrm{in}}$ makes the source vanish exactly at the inner
boundary. In \code{"outer"} mode only the first role moves; $C$ keeps the
inner slope from \code{dne_dx_inner} / \code{bc_origin}, because it is a
model constant rather than a boundary condition.

\paragraph{Implementation.} \code{solve_sc_scipy} imposes the
\code{"outer"} pair directly (\code{solve_bvp} accepts both residuals at one
end). \code{solve_sc_firedrake} cannot: the strong Dirichlet condition at
boundary id 2 makes $w(0)=0$, so a \code{ds(2)} flux term vanishes
identically. There the equivalent well-posed statement --- the inner flux
is an unknown fixed by the separatrix gradient --- is solved by secant
iteration on the \code{ds(1)} slope (\code{_shoot_outer_grad_sc},
Section~\ref{sec:shoot}), leaving the function space and the Picard loop
untouched. Diagnostics land in \code{self.grad_bc_info}.

%=====================================================================
\section{Class \code{OneDSolverMixin}}
%=====================================================================
Original (implicit) Saarelma--Connor model, non-dimensionalised (``1D'').
A mixin over \code{SaarelmaConnorBase}; it is combined with the base and
\code{ThreeDSolverMixin} into the single public
\code{saarelma_connor}. It solves the single second-order ODE for $n_e$
only (report Eqs.~(6)--(7) / Saarelma Eq.~16 and corrected Eq.~15).

After a successful solve the following attributes are set:

\begin{center}
\begin{tabular}{ll}
\toprule
Attribute & Meaning \\
\midrule
\code{x_sol}, \code{ne_sol}, \code{dne_dx_sol} & SI profiles \\
\code{hat_x_sol}, \code{hat_ne_sol} & non-dimensional profiles \\
\code{E_sol} & converged FC kernel $E(x)$ on \code{x_sol} \\
\code{alpha_bar_ped}, \code{kbm_gate_on} & final KBM gate state \\
\code{picard_info}, \code{kbm_info} & iteration diagnostics \\
\code{_L_sc}, \code{_n0_sc}, \code{_S0_sc}, \code{_V0_sc}, \code{_D0_sc} & reference scales \\
\bottomrule
\end{tabular}
\end{center}

%=====================================================================
\section{Method reference}
%=====================================================================

\subsection{\code{_check_free_params_sc()}}
Backwards-compatible alias for \code{check_free_params}.

\subsection{\code{_ensure_sc_setup(x_res, free_params=None, force=False)}}
Equilibrium-only setup shared by both implementations. Optionally updates
the free parameters, then runs the parent initialisation methods (solver
grids, form factors, ETG coefficient), locates the inner boundary, and
computes the non-dimensional reference scales
$L=|x_{\mathrm{inner}}|$, $n_0=\code{ne_x0}$, $S_0=S_i(0)$,
$[V]_0 = L n_0 S_0$, $[D]_0 = L^2 n_0 S_0$.

\begin{description}[style=nextline]
  \item[\code{x_res} (int)] Resolution handed to \code{setup_solver_grids}.
  \item[\code{free_params} (dict or None)] Optional
  \code{{'alpha_crit', 'De_chie_etg', 'C_KBM', 'nFC_x0'}} (plus any keyword
  accepted by \code{update_free_params}, e.g.\ \code{psi_N_inner_boundary})
  applied before the setup.
  \item[\code{force} (bool)] Re-run \code{setup_solver_grids} /
  \code{form_factor} even if they were already built for this \code{x_res}.
\end{description}
If \code{psi_N_inner_boundary} is \code{None}, \code{find_inner_boundary}
is called; it needs a $D_{\mathrm{KBM}}$ estimate, and a zeroed array (KBM
off) is used as the conservative pre-solve choice. Raises if
$x_{\mathrm{inner}}\ge 0$ or if $S_i(0)$ is not a usable (finite,
positive) rate.

\subsection{\code{_shoot_outer_grad_sc(F, N, bcs, snes_params, slope_c, target, tol, max_it, verbose, tag)}}
\label{sec:shoot}
Solve $F=0$ subject to $N'(0)=\code{target}$.

The inner boundary carries no condition of its own in the \code{"outer"}
mode, so the \code{ds(1)} flux slope \code{slope_c} is the free unknown; a
damped secant iteration drives
\[
r(s) = N'(0;s) - \code{target}
\]
to zero (with tolerance \code{tol} relative to $|\code{target}|$).
\code{slope_c} is deliberately \emph{not} the constant $C$ of the source
term --- that stays at the inner slope --- so only the boundary flux moves
here.

\paragraph{Step control.} The seed is the inner slope, which can be far from
the answer (the free flux routinely lands on the other side of zero), and
Eq.~(7) need not be solvable for every slope in between: each step is
capped at four times the previous accepted one and backtracked, halving (up
to 12 times), whenever SNES fails. The very first step is bootstrapped
with a finite-difference step of $0.05|s|$ (or $10^{-3}$ if $s=0$).

\paragraph{Fixed restart state.} Every attempt restarts from the
\emph{same} reference state (the iterate this routine was handed) rather
than from the previous attempt's answer. That is deliberate. Chaining warm
starts makes $r(s)$ path dependent --- revisiting one slope reproduced
$N'(0)$ only to $\sim5\times10^{-5}$ on the DIII-D 158091 case --- and,
worse, once two slopes are close the warm start already passes SNES's
convergence test, so it returns at iteration zero with the density
untouched, and the secant sees a \emph{bit-identical} residual for two
different slopes and concludes the gradient does not respond at all.
Restarting from a fixed reference makes $r(s)$ a genuine function of $s$,
which is what the secant needs; the cost is a few more Newton steps per
attempt. For the same reason \code{snes_stol} is set to 0 so SNES cannot
exit on the step-size test before it has done any work.

\paragraph{Returns} a diagnostics dict with keys \code{converged},
\code{iterations}, \code{solves}, \code{backtracks}, \code{slope},
\code{rel_residual}, \code{history}. Raises \code{RuntimeError} if the
seed solve fails, if the secant stalls (zero denominator), if SNES keeps
failing through all backtracks, or if \code{max_it} is exhausted.

\subsection{\code{calc_D_KBM_average_sc(n_e_ped, x_ped)}}
\label{sec:kbm}
Pedestal-averaged-$\alpha$ KBM diffusivity on the \code{x_init} grid. This
is the ``average'' Picard gate mode of the 3D solver's
\code{calc_pressure_quantities_nondim}, specialised to the original model.
The Connor--Hastie $\alpha$ is evaluated locally from the given pedestal
density profile,
\[
\alpha(x) = \alpha_{\mathrm{nodp}}(x)\,\frac{d}{dx}\bigl[n_e(T_e+T_i)e\bigr],
\]
(neutral pressures neglected, quasi-neutral plasma assumed), averaged over
the pedestal, and gated as a whole:
\[
\text{gate} = \bar\alpha > \alpha_{\mathrm{crit}},\qquad
D_{\mathrm{KBM}} = (\bar\alpha - \alpha_{\mathrm{crit}})\,
   \frac{C_{\mathrm{KBM}}\,c_s\,\rho_s^2}{a}
\]
everywhere when the gate is on, zero otherwise.

\begin{description}[style=nextline]
  \item[\code{n_e_ped} (array)] Electron density (m$^{-3}$) on \code{x_ped}.
  \item[\code{x_ped} (array)] Ascending pedestal grid (m), $x\in[x_{\mathrm{inner}},0]$.
  \item[Returns] \code{D_KBM_xinit}: KBM diffusivity (m$^2$/s) on
  \code{self.x_init} (used to build the coefficient interpolators /
  Functions).
  \item[Sets] \code{alpha_bar_ped} (float), \code{alpha_frac_above}
  (float), \code{kbm_gate_on} (bool), \code{alpha_local_ped} (local
  $\alpha$ on \code{x_ped}), \code{pres_ped}, and \code{_D_KBM} (same as
  the return value; also keeps \code{find_inner_boundary} re-entrant).
\end{description}

\subsection{\code{_exp_kernel_sc(n_e, x)}}
FC attenuation kernel (exponent of Saarelma Eq.~11)
\[
E(x) = \exp\!\left[\int_0^x \frac{n_e(S_i+\SCX)}{\fFC\VFC}\,dx'\right].
\]
\code{x} is an ascending grid ($x_{\mathrm{inner}}\dots0$). Returns $E$ on
the same grid; $E(0)=1$ and $E$ decays monotonically inward. This is the
same closure the parent class uses for $\nFC$ in
\code{find_inner_boundary} / \code{compute_post_solve_SC_neutrals}, i.e.\
$\nFC(x) = \code{nFC_x0}\cdot E(x)$. Because $x<0$, the cumulative
integral is taken on a descending ordering (separatrix $\to$ core) so it
starts at 0; the result is negative since the integrand is positive.

\subsection{\code{_C_cx_sc(x)}}
CX-flux coefficient $\CCX(x)$ of report Eq.~(7), Eq.~\eqref{eq:Ccx},
evaluated on \code{x} (m). It is purely a ratio of velocities and rate
coefficients, so it is scale-invariant and built once in SI.

\subsection{\code{_post_solve_neutrals_sc()}}
$\nFC$ and $\nCX$ (m$^{-3}$) on \code{self.x_sol} from the converged $n_e$.
The neutrals are not unknowns of this model, so they are reconstructed
after the solve with Saarelma \emph{et al.} (2023) Eqs.~(11)--(12):
\begin{align*}
\nFC &= \code{nFC_x0}\cdot E(x),\\
\nCX &= \max\!\left(-\frac{\gradr\Dped}{\VCX\fCX}\bigl(n_e' - n_e'(-\infty)\bigr)
       + (\CCX-1)\nFC,\; 0\right).
\end{align*}
This is the same closure as the base class's
\code{compute_post_solve_SC_neutrals}, with two differences that make it
match what this solver actually solved:
\begin{itemize}[nosep]
  \item every coefficient is interpolated from \code{x_init} onto
  \code{x_sol}, rather than assumed to share its length --- the adaptive
  \code{solve_bvp} mesh never does;
  \item the ETG part of $\Dped$ is $C_{\mathrm{ETG}}/n_e$, which carries
  \code{De_chie_etg}, rather than $D_{\mathrm{ETG},x}/n_e$, which does not.
  The two agree only for \code{De_chie_etg} $=1$.
\end{itemize}
$D_{\mathrm{KBM}}$ is the frozen Picard value on \code{self._D_KBM}, i.e.\
the one the final $n_e$ was solved with.

\subsection{\code{_check_eq6_form_sc(eq6_form)} (static)}
Validate and normalise the \code{eq6_form} switch (\code{"complete"} or
\code{"paper"}). Choosing \code{"paper"} emits a \code{RuntimeWarning}:
it drops the shape factor $1+\SCX/(2S_i)$ and cancels $\gradr$ against
$\Dped$; on DIII-D 158091 this makes it $\sim35\times$ stiffer than the
complete form and no implementation can solve it in double precision.

\subsection{\code{_check_first_step_sc(first_step)} (static)}
Validate and normalise the \code{first_step} switch (\code{"auto"},
\code{"eq6"}, \code{"skip"}).

\subsection{\code{_first_step_failed_sc(tag, first_step, err)}}
Handle a failed step 1: re-raise (\code{first_step="eq6"}), or warn and
fall back to \code{"skip"} (\code{"auto"}), setting
\code{first_step_used = "skip (Eq. 6 failed)"}.

Eq.~(6) is a positive-feedback amplifier --- every bit of extra density
gradient increases the implied $\nFC$, which steepens the gradient further
--- so for a steep enough inner-boundary slope it has no solution that
still reaches $n_e(0)=\code{ne_x0}$ (the profile collapses first). The
amplification is $\exp(\int C_{A6} N\,d\xi)$, controlled by the
inner-boundary gradient, the pedestal width and the $\Dped$ shape.
Saarelma \emph{et al.} (2023) only use it as ``a convenient starting
point'' for the Eq.~(7) iteration, so when it fails the Picard loop can
equally well start from the initial guess. Suggested remedies:
\code{first_step="skip"}, a shallower \code{dne_dx_inner}, or an inner
boundary closer to the separatrix.

\subsection{\code{_record_picard_sc(tag, converged, n_it, history, picard_max_it, picard_relax, picard_rtol)}}
Store \code{picard_info} / \code{kbm_info} and raise if the Picard loop
diverged. The usual culprit for non-convergence is a KBM gate limit cycle:
$\bar\alpha$ straddles $\alpha_{\mathrm{crit}}$, so the gate (and with it
$D_{\mathrm{KBM}}$) flips every iteration and the profile alternates
between two states. This is diagnosed explicitly from the last six
iterations (the fix --- \code{picard_relax} $<1$, or moving
$\alpha_{\mathrm{crit}}$ away from $\bar\alpha$ --- is different from
``just iterate longer'').

%---------------------------------------------------------------------
\subsection{\code{solve_sc_scipy(...)}}
\label{sec:scipy}
%---------------------------------------------------------------------
Implicit scipy (\code{solve_bvp}) solver for the original model. Step 1
solves the non-dimensional no-CX equation (report Eq.~6) for a first $n_e$;
step 2 Picard-iterates the full equation (report Eq.~7) with the
exponential kernel $E(x)$ and the pedestal-averaged KBM diffusivity frozen
from the previous iterate, until the profile and the KBM gate stop
changing.

\paragraph{Parameters.}
\begin{description}[style=nextline]
  \item[\code{x_res} (int, 200)] Number of points of the (uniform)
  non-dimensional solver grid handed to \code{solve_bvp} as its initial
  mesh.
  \item[\code{free_params} (dict or None)] Optional free parameters applied
  before solving; see \code{_ensure_sc_setup}.
  \item[\code{bc_origin} (\code{"p-file"} | \code{"user"})] Where the
  prescribed slopes come from: the p-file density gradient, or the
  user-supplied \code{dne_dx_inner} / \code{dne_dx_outer}.
  \item[\code{dne_dx_inner} (float or None)] SI inner-boundary slope
  (m$^{-4}$, must be negative) when \code{bc_origin="user"}. Always supplies
  Saarelma's constant of integration $C$ in the source term, whatever
  \code{ne_grad_bc_loc} is; it is additionally the Neumann BC value in
  \code{ne_grad_bc_loc="inner"} mode.
  \item[\code{ne_grad_bc_loc} (\code{"inner"} | \code{"outer"})] Which end
  carries the prescribed $dn_e/dx$; see Section~\ref{sec:bcs}.
  \code{"outer"} puts both $n_e$ conditions at the separatrix and leaves the
  inner boundary free.
  \item[\code{dne_dx_outer} (float or None)] SI separatrix slope
  $dn_e/dx|_{x=0}$ (m$^{-4}$, must be negative) used when
  \code{ne_grad_bc_loc="outer"}. Given explicitly it wins for any
  \code{bc_origin} (this is how Saarelma Eq.~(20) is fed in); left
  \code{None} it is read off the p-file gradient at $x=0$, which requires
  \code{bc_origin="p-file"}.
  \item[\code{dne_dx_neginf} (float or None)] Override for the integration
  constant $C$ (resolved by \code{bc_ig_helpers.resolve_integration_constant}).
  \item[\code{initial_guess} (\code{"pfile"} | \code{"linear"} | \code{"tanh"})]
  Shape of the initial $n_e$ profile (SI).
  \item[\code{tanh_width}, \code{tanh_center} (float or None)] Parameters
  of the \code{"tanh"} initial guess (SI metres).
  \item[\code{eq6_form} (\code{"complete"} | \code{"paper"})] Which form of
  the first step to solve. \code{"paper"} is Eq.~6 exactly as printed and
  is only usable when $\gradr$ is close to 1: otherwise it either fails
  outright or converges to a collapsed profile that the Eq.~(7) iteration
  cannot then start from.
  \item[\code{first_step} (\code{"auto"} | \code{"eq6"} | \code{"skip"})]
  Whether to run step 1 at all. \code{"eq6"} always runs it and raises if it
  fails; \code{"skip"} starts the Eq.~(7) Picard loop directly from
  \code{initial_guess}; \code{"auto"} (default) tries Eq.~(6) and falls
  back to \code{"skip"} with a \code{RuntimeWarning} if it has no solution.
  \item[\code{picard_max_it} (int, 50)] Maximum number of Picard iterations
  of the full Eq.~(7).
  \item[\code{picard_rtol} (float, $10^{-6}$)] Relative max-norm tolerance
  on the change in $n_e$ between Picard iterations.
  \item[{\code{picard_relax} (float in $(0,1]$)}] Under-relaxation of the
  frozen $D_{\mathrm{KBM}}$ and $E(x)$ updates.
  \item[\code{bvp_tol} (float)] \code{solve_bvp} residual tolerance.
  \item[\code{bvp_max_nodes} (int)] \code{solve_bvp} maximum number of mesh
  nodes.
  \item[\code{reuse_setup} (bool)] Reuse cached solver grids / form
  factors when the resolution is unchanged (set \code{False} after
  changing the equilibrium).
  \item[\code{verbose} (bool or None)] Override \code{self.verbose}.
\end{description}

\paragraph{Returns.} The common solver result schema (see
\code{SaarelmaConnorBase._build_result_dict}). \code{x}, \code{ne} and
\code{dne_dx} are the converged SI profiles; \code{nFC} / \code{nCX} are
derived from them via Saarelma Eqs.~(11)--(12), since they are not unknowns
of this model.

\paragraph{Internal helper \code{_build_f_interp(D_KBM_xinit)}.}
Returns interpolators for $f_0$, $f_1$ and their $x$-derivatives, with
$f_0 = \gradr(D_{\mathrm{NEO}}+D_{\mathrm{KBM}})$ (m$^2$/s) and
$f_1 = \gradr C_{\mathrm{ETG}}$ (m$^2$/s $\cdot$ m$^{-3}$). The expanded
ODE needs $df_0/dx$ and $df_1/dx$ (the $C_K$ term), and the accuracy of the
whole solve is limited by them: \code{x_init} is the p-file pressure grid,
which typically carries only a few tens of points across the pedestal. A
shape-preserving (PCHIP) interpolant of $f_0$ / $f_1$ with its analytic
derivative is markedly better there than \code{np.gradient} followed by
linear interpolation, and unlike a natural cubic spline it cannot overshoot
on $C_{\mathrm{ETG}}$, which spans several decades over \code{x_init}. The
conservative Firedrake form never differentiates $f$, so it has no
equivalent of this term.

\paragraph{Internal helper \code{_conductance(xi, N, ...)}.}
Returns $x$, $f$, $C_K$, $C_N$ at $(\xi,N)$. $N$ is floored at $10^{-8}$
while it appears in a denominator, so that a solver excursion towards
$N\le0$ cannot produce a division by zero.

\paragraph{Boundary residuals.} In \code{"inner"} mode:
$Y_a[1]-N'_{bc}$ (Neumann at $\xi=-1$) and $Y_b[0]-1$ (Dirichlet at
$\xi=0$). In \code{"outer"} mode both residuals sit at $\xi=0$:
$Y_b[0]-1$ and $Y_b[1]-N'_{bc}$. \code{solve_bvp} only needs the right
\emph{number} of residuals, so this needs no shooting.

\paragraph{Step 1 details.} $D_{\mathrm{KBM}}$ is frozen from the initial
guess. A ``converged'' step-1 profile with $\min N\le0$ is rejected: Eq.~6
can be satisfied by a density that crosses zero inside the pedestal, which
no Eq.~7 iterate can recover from (and which is unphysical anyway). If step
1 is skipped or fails, the Picard loop starts from the initial guess on the
uniform grid.

\paragraph{Step 2 details.} Each iteration refreezes $E(x)$ and
$D_{\mathrm{KBM}}$ (averaged-$\alpha$ gate) from the previous iterate with
optional under-relaxation, warm-starts \code{solve_bvp} from the previous
iterate on the uniform grid, and checks convergence on the relative
max-norm change of $n_e$ together with a stable gate state.

%---------------------------------------------------------------------
\subsection{\code{_ensure_firedrake_mesh_sc(mesh_n, fe_degree, force=False)}}
%---------------------------------------------------------------------
Build (or reuse) the non-dimensional mesh on $\xi\in[-1,0]$ (CG elements of
degree \code{fe_degree}). Only the mesh, function space and DOF coordinates
are cached: the coefficient Functions depend on the free parameters and the
reference scales, so they are rebuilt on every solve by
\code{_build_sc_coefficients}. Returns \code{(mesh, V, xi_dofs)}.

\subsection{\code{_build_sc_coefficients(V, x_dofs_si)}}
Frozen equilibrium coefficient Functions in non-dimensional units. All
entries are \code{Function}s on \code{V} in DOF order:

\begin{center}
\begin{tabular}{ll}
\toprule
Key & Quantity \\
\midrule
\code{g} & $\gradr$ (dimensionless) \\
\code{hat_Si}, \code{hat_Scx} & $S_i/S_0$, $\SCX/S_0$ \\
\code{hat_Vcx} & $\VCX/[V]_0$ \\
\code{fFC}, \code{fCX} & form factors (dimensionless) \\
\code{hat_D_NEO} & $D_{\mathrm{NEO}}/[D]_0$ \\
\code{hat_C_ETG} & $C_{\mathrm{ETG}}/(n_0[D]_0)$ (so $\hat C_{\mathrm{ETG}}/N$ is the non-dim ETG diffusivity) \\
\code{C_cx} & report Eq.~(7) CX-flux coefficient (already dimensionless) \\
\bottomrule
\end{tabular}
\end{center}

%---------------------------------------------------------------------
\subsection{\code{solve_sc_firedrake(...)}}
\label{sec:firedrake}
%---------------------------------------------------------------------
Implicit Firedrake solver for the original model. Solves the same
non-dimensional two-step problem as \code{solve_sc_scipy} but in
conservative (weak) form with CG finite elements and SNES Newton for each
frozen-coefficient nonlinear solve:

\paragraph{Step 1} (report Eq.~6, no CX; \code{eq6_form="paper"}):
\[
F_6 = \int \hat f N' w'\,d\xi
 + \int N\,\frac{\hat f}{g}\,\frac{(\hat S_i+\hat S_{\mathrm{CX}})(N'-N'_{\mathrm{in}})}{\hat V_{\mathrm{FC}}\fFC}\,w\,d\xi
 + \hat f N'_{\mathrm{in}}\,w\,ds(1) = 0.
\]
\code{eq6_form="complete"} replaces $\hat f/g$ by
$\hat f/(1+\hat S_{\mathrm{CX}}/(2\hat S_i))$.

\paragraph{Step 2} (report Eq.~7, Picard on $E(x)$ and $D_{\mathrm{KBM}}$):
\[
F_7 = \int \hat f N' w'\,d\xi
 + \int N\hat S_i\left[\frac{\hat f(N'-N'_{\mathrm{in}})}{\hat V_{\mathrm{CX}}\fCX}
   - \CCX\,\hat n_{\mathrm{FC},0}\,E\right] w\,d\xi
 + \hat f N'_{\mathrm{in}}\,w\,ds(1) = 0,
\]
with $\hat f = g(\hat D_{\mathrm{NEO}}+\hat D_{\mathrm{KBM}}) + g\hat C_{\mathrm{ETG}}/N$
(the ETG piece keeps its $1/N$ dependence inline, so Newton differentiates
it exactly), Dirichlet $N(0)=1$ (boundary id 2), and the Neumann inner flux
imposed weakly through \code{ds(1)} (boundary id 1).

The frozen KBM diffusivity \code{hat_D_KBM} and kernel $E$ are Functions
refreshed between Picard iterations from the latest density --- the same
``average''-gate Picard pedestal-pressure treatment as the 3D solver
(Saarelma Eqs.~24--25 diffusivity frozen into the $D$-slot, optional
under-relaxation). $E$ is unused by step 1 but seeded from the initial
guess so the form is well defined from the start.

\paragraph{\code{"outer"} mode.} With \code{ne_grad_bc_loc="outer"} the
Dirichlet $N(0)=1$ stays but the inner boundary is freed and its
\code{ds(1)} flux becomes the unknown that a secant iteration adjusts until
$N'(0)$ matches \code{dne_dx_outer} (see Section~\ref{sec:bcs} for why a
\code{ds(2)} term cannot impose it directly). Every nonlinear solve ---
step 1 and each Picard iteration of step 2 --- is wrapped in that
iteration, with the slope carrying over so later shootings start close.
The source constant $N'_{\mathrm{in}}$ is untouched by this and keeps the
inner slope. Two constants are therefore used: \code{dN_in_c} (the source
constant $C$, always the inner slope) and \code{dN_bc_c} (the \code{ds(1)}
flux: the BC value in \code{"inner"} mode, the secant unknown in
\code{"outer"} mode).

\paragraph{Secant seed.} In \code{"outer"} mode the secant is seeded with
zero inner flux (or \code{grad_bc_seed} if given). Zero inner flux is the
natural boundary condition a Galerkin form defaults to, so it is the
neutral starting point and --- unlike seeding from a pedestal-top gradient
--- uses no information the \code{"outer"} pathway forbids. It is also
markedly more robust: $r(s)$ is discontinuous where the profile switches
solution branch, and approaching from zero avoids straddling that jump.

\paragraph{Step-1 failure.} If step 1 fails (or converges to a non-positive
profile), $N$ is restored to the initial guess and \code{dN_bc_c} to its
seed: a diverged SNES leaves $N$ in whatever state the failed line search
reached, and in \code{"outer"} mode the secant may have walked
\code{dN_bc_c} somewhere unhelpful for the Eq.~7 restart.

\paragraph{Parameters} are as in \code{solve_sc_scipy}, plus the finite
element / PETSc controls \code{fe_degree}, \code{linear_solver},
\code{ksp_rtol}, \code{ksp_max_it}, and the \code{"outer"}-mode secant
controls \code{grad_bc_tol} (relative tolerance on $N'(0)$, default
$10^{-8}$), \code{grad_bc_max_it} (default 25) and \code{grad_bc_seed}
(SI seed slope, default zero flux). Note that the Firedrake
\code{picard_rtol} (default $10^{-8}$) is a relative $L^2$ tolerance.

\paragraph{Returns.} The common solver result schema; profiles are on the
sorted DOF grid. In \code{"outer"} mode the inner slope selected by the
separatrix-gradient condition is reported as
\code{self.dne_dx_inner_solved} (distinct from \code{self.dne_dx_inner},
which is the source constant $C$ and was an input).

\subsection{\code{solve_sc(solver_structure="firedrake", **kwargs)}}
Solve the original model with the chosen solver structure
(\code{"firedrake"} or \code{"scipy"}); \code{**kwargs} are passed
straight through to \code{solve_sc_firedrake} or \code{solve_sc_scipy}.
Returns the common solver result schema.

\subsection{Module \code{__getattr__(name)}}
Resolves the pre-refactor class name \code{saarelma_connor_sc} to the
assembled class \code{saarelma_connor}. Deferred rather than a top-level
import: \code{saarelma_connor_api} imports this module, so binding the name
at import time would be a cycle.

\end{document}
